\section{Mistral、Voxtral 3と音声エージェント基盤を発表}
MistralはVoxtral 3を公式Xで発表した。多言語の音声文字起こしを担うSmall、Mini、Realtimeなどに加え、音声生成モデルとLiveKitベースのvoice agent SDKを提供するという。API利用と自己ホストの両方に対応すると説明する。低遅延の発話終了検出や、35言語以上で感情・話速を調整できる音声生成も訴求しているが、これらの性能評価は発表者の説明に基づく。今後は実環境での遅延や品質の比較が焦点となる。
\dailyxsource{Mistral AI (@MistralAI)}{https://x.com/MistralAI/status/2112574615318876178}

\section{DeepSeek、推論基盤DeepInferを公開}
DeepSeekが推論フレームワーク「DeepInfer」をApache-2.0ライセンスで公開したとの情報が広がった。Grokの記録では、DeepSeek-V4.1-Flash、NVIDIA SM100、Huawei Ascend 950への対応が説明される一方、KVキャッシュ関連など一部の構成要素は公開対象に含まれないという。生産環境で使われているとの主張や高スループット・低遅延の評価もあるが、再現可能な比較結果は本号では確認していない。公式X投稿の直接URLも未特定のため、推測したリンクは付さない。

\section{Claudeによるサソリ毒素結合タンパク質の設計}
Anthropicは、内部研究モデルのClaudeを使ってサソリ神経毒に結合するタンパク質を設計した実験を紹介した。約64時間の設計作業を経て、試験した211候補のうち17候補が実験室で結合基準を満たしたと報告する。設計候補は30万件以上に上るという。これは結合の確認であり、毒性を中和する抗毒治療が成立したことを意味しない。同社は研究の初期段階として位置づけており、実用化には中和効果や安全性の検証が必要となる。
\dailyxsource{Anthropic (@AnthropicAI)}{https://x.com/AnthropicAI/status/2112579251555721573}

\section{OpenAI、生命科学企業Orchestra Bioを買収}
OpenAIはOrchestra Bioの買収を公式Xで発表した。同社はバイオ医薬品の開発プロセスを調整するプラットフォームを手掛けており、OpenAIは生命科学向けAIモデル・ツール「Rosalind」との連携によって治療開発を速める方針を示した。AIの研究支援を、実際の創薬・開発工程につなげる取り組みとして注目される。一方、買収金額や統合後の製品計画について、今回のGrok収集では追加の詳細が示されていない。
\dailyxsource{OpenAI (@OpenAI)}{https://x.com/OpenAI/status/2112607347688583586}

\section{Grok Botのファイル対応拡大などが話題に}
Elon Musk氏がGrok Botの大きな更新をXで強調し、関連するデモや機能紹介が共有された。コミュニティ投稿では、送信ファイルの上限が250MB、動画が1GB、同時に最大20ファイルに拡大したほか、通知やフォームの事前入力が改善されたとされる。ただし、個々の数値や機能一覧はGrokの収集で公式リリースノートまで照合できていない。更新の存在に関する投稿と、具体的な仕様の確定を分けて読む必要がある。
\dailyxsource{Elon Musk (@elonmusk)、関連投稿}{https://x.com/elonmusk/status/2112589019380875714}

\section{GeminiのPersonal Intelligence、Android Autoへ}
Googleは、GeminiのPersonal IntelligenceをAndroid Autoへ展開し始めたと公式Xで告知した。Gmailやカレンダーなど、利用者が許可したアプリの情報を組み合わせ、運転中の複雑な依頼にもハンズフリーで応える体験を目指す。将来的にはGoogle built-in搭載車への対応も計画されているという。車載AIの個人向け機能が具体化する一方、利用可能な地域・端末やデータ連携の制御については、追加の案内を待つ必要がある。
\dailyxsource{Google (@Google)}{https://x.com/Google/status/2112655614132285592}

\section{Reflection Beam 1.0、Vals Indexでの評価を報告}
Vals AIは、Reflectionのオープンウェイトモデル「Beam 1.0」の評価結果をXで紹介した。投稿によれば、Indexでのタスク実行時間は平均約7分で、オープンウェイトモデルの中でも高速とされる。ランキング上位、長いコンテキストへの対応、無料公開なども主張されている。ただし、順位や速度は評価側の発表に基づき、他のベンチマークで同じ傾向が出るかは未確認である。実利用での性能と公開条件の確認が次の論点となる。
\dailyxsource{Vals AI (@ValsAI)}{https://x.com/ValsAI/status/2112655199999590772}

\section{DeepMind、Gemini 4 Argonのデモを公開}
Google DeepMindは、Gemini 4 Argonを用いたデモを公式Xで紹介した。楽曲に同期して動く3D周期表や、ブラウザ上でのリアルタイムな毛皮の物理表現など、対話から視覚的・動的な成果物を作る能力を示す内容である。Antigravity上での制作にも言及された。これらはデモの紹介であり、モデルの一般提供範囲や仕様、他条件での再現性を保証するものではない。正式な利用条件については引き続き確認が必要だ。
\dailyxsource{Google DeepMind (@GoogleDeepMind)}{https://x.com/GoogleDeepMind/status/2112557925264752930}

\section{Vercel AI SDKにBash Toolが登場}
VercelのAI SDKは、エージェントにシェル操作を提供する「Bash Tool」の導入を公式アカウントで発表した。任意のモデルとAI Gatewayを組み合わせて使えると説明している。コード編集や開発作業を支援するエージェントにとって、シェル実行は重要な機能となる。一方、コマンド実行に伴う権限管理やサンドボックスの制約、具体的な利用例については今回のGrok収集では確認されていない。
\dailyxsource{AI SDK (@aisdk)}{https://x.com/aisdk/status/2112662275660349720}

\section{llm-d 1.0、分散推論スタックとして紹介}
Google Cloud上のTPU・GPU向け分散推論スタック「llm-d 1.0」が公開されたとの投稿が共有された。オープンソースの構成を通じて、エージェント向け推論のコストやスループットを改善するという。Grokの記録にはGoogle Cloudのブログへのリンクがあるが、今回の収集では実装の詳細や性能比較までは確認していない。特定クラウドの環境以外への移植性や、実際の運用条件での効果が今後の注目点となる。
\dailyxsource{GCP Weekly (@gcpweekly)、関連情報}{https://x.com/gcpweekly/status/2112506828428918816}

\section{Grok Botによる障害対応自動化の報告}
SpaceXSIがGrok Botをサポートチケットやインシデント対応に活用したとする投稿が共有された。投稿では、204件のチケットを約10分で処理したほか、複数のボットを組み合わせた運用が紹介されている。ただし、Grokは元のブログなどの一次資料を確認できておらず、件数や処理時間、実際の運用成果は未検証である。エージェントによる業務自動化の事例として話題になったものの、効果を確定的に評価する段階にはない。
\dailyxsource{X Freeze (@XFreeze)、未検証の事例紹介}{https://x.com/XFreeze/status/2112624561737158831}

\section{Claude Code 2.1.305の更新情報が共有}
非公式の変更ログアカウントが、Claude Codeのバージョン2.1.305に関する更新を報告した。Bashコマンド実行、サブエージェント、サンドボックスなどの改善が挙げられているが、今回のGrok収集では公式リリースノートとの照合はできていない。開発エージェントの実行環境や安全性に関わる変更として関心を集めた一方、各機能の正式な提供状況や変更内容は未確認の情報として扱う。
\dailyxsource{Claude Code Changelog (@ClaudeCodeLog)、非公式}{https://x.com/ClaudeCodeLog/status/2112650002778501148}
